Семейства обобщённых последовательностей Фибоначчи, зависящие
от параметров, рассматриваются, например, в~\cite{Mehraban2006Gulliver}.
Чтобы однозначно зафиксировать индексацию, в настоящей работе
используется следующая нормировка. Она выбрана так, чтобы при $p=1$
порядок рекурсии был равен $k$.

\begin{definition}\label{defn1pil}
Пусть $k\ge2$, $p\ge1$, $t\ge1$ --- целые числа и $r=k+p-1$.
Обобщённая $(t,k)$-$p$-последовательность
$(F_n^p(t,k))_{n\ge0}$ определяется равенствами
\begin{equation}
 F_n^p(t,k)=tF_{n-1}^p(t,k)
       +\sum_{j=p+1}^{r}F_{n-j}^p(t,k),
 \qquad n\ge r,
\end{equation}
и начальными условиями
\begin{equation}
 F_0^p(t,k)=\cdots=F_{r-2}^p(t,k)=0,
 \qquad F_{r-1}^p(t,k)=1.
\end{equation}
При необходимости полагаем $F_n^p(t,k)=0$ для $n<0$.
\end{definition}

Начальные условия задают ровно $r$ первых членов, а рекурсия
определяет каждый следующий член через ранее заданные.
По индукции последовательность существует и единственна.
В частности, при $p=1$ начальные условия имеют вид
$F_0^1=\cdots=F_{k-2}^1=0$, $F_{k-1}^1=1$, и
\[
 F_n^1(t,k)=tF_{n-1}^1(t,k)+F_{n-2}^1(t,k)+\cdots+F_{n-k}^1(t,k)
 \qquad(n\ge k).
\]
При $t=1$, $k=2$, $p=1$ получается обычная последовательность
Фибоначчи с $F_0=0$, $F_1=1$.
